\chapter{Dual complexes and log discrepancy over the complex numbers}
\index{dual complex!complex degeneration}\index{log discrepancy!complex}
\label{chap:mp-dual-complexes}

The valuative space attached to a compact complex manifold is not built from
its field of meromorphic functions.  That field can be too small: a very
general complex torus has algebraic dimension zero.  Instead, one evaluates
all coherent ideals.  Modifications then recover the familiar divisorial and
monomial points, while equivariant modifications of
\(X\times\mathbf P^1\) organize those points into dual complexes.

This chapter develops the birational geometry needed in
\Cref{chap:mp-psh,chap:mpwn-cy,chap:mpwn-k-stability}.  The primary source is
Sections~1--3 of \cite{MP24}.  One input in that source deserves to be made
visible: Bingener's relative analytification over Stein algebras is used to
compare analytic divisors with divisors over analytic local rings.  That
comparison is indispensable in the proof of divisorial density, and density
is in turn used in the proof of the inverse-limit description.  We therefore
give the three arguments in that order and cite the original source
\cite{Bingener76Stein}; the detailed reconstruction in
\cite[Chapter~5]{XiaStein26} supplies the modern proof.

Throughout, \(X\) is a connected compact K\"ahler manifold unless a normal
compact K\"ahler space is explicitly allowed.  The tropical-spectrum and
Stein-local statements below remain meaningful more generally, but the
cofinal system of test configurations used for divisorial density and dual
complexes is the K\"ahler system of \cite[Definition~2.5.1]{MP24}; we do not
assert those model statements for an arbitrary non-K\"ahler compact complex
space.  We write \(\mathfrak I_X\) for the semiring of coherent ideal sheaves
on \(X\), and \(X^{\bethan}\) for its tropical spectrum.

The dependency chain should be kept in view.  Coherent ideals first produce
continuous tropical coordinates.  Normalized blowups turn those coordinates
into vertical Cartier data.  Bingener's comparison then globalizes local
divisorial centers, which yields divisorial density.  Only after density is
available can the compatible retractions to dual complexes be proved to
remember every point.  In particular, the inverse-limit theorem is not a
formal consequence of resolution of singularities alone.

\section{Coherent ideals and analytic modifications}

A point \(v\in X^{\bethan}\) is a map

\[
 v:\mathfrak I_X\longrightarrow[0,+\infty]
\tag{59.1}\label{eq:mp-semivaluation-ideals}
\]

satisfying

\[
 \begin{split}
 v(\mathfrak a\mathfrak b)&=v(\mathfrak a)+v(\mathfrak b),\\
 v(\mathfrak a+\mathfrak b)&=
             \min\{v(\mathfrak a),v(\mathfrak b)\},\\
 v(\cO_X)&=0,\qquad v(0)=+\infty.
 \end{split}
\tag{59.2}\label{eq:mp-semivaluation-axioms}
\]

The topology is pointwise convergence on coherent ideals.

\begin{proposition}\label{prop:mp-tropical-compactness}
The space \(X^{\bethan}\) is compact and Hausdorff.
\end{proposition}

\begin{proof}
Give \([0,+\infty]\) its compact order topology and consider
\[
                    [0,+\infty]^{\mathfrak I_X}.
\]
It is compact by Tychonoff's theorem and Hausdorff.  Extended addition on
\([0,+\infty]\), with \(a+(+\infty)=+\infty\), and the minimum operation
are continuous.  Therefore each identity in
\eqref{eq:mp-semivaluation-axioms}, for each fixed pair of coherent ideals,
cuts out a closed subset of the product.  The normalizations
\(v(\cO_X)=0\) and \(v(0)=+\infty\) are closed as well.  Their intersection
is exactly \(X^{\bethan}\), so it is compact.  Coordinate evaluations
separate distinct semivaluations, giving Hausdorffness.
\end{proof}

If \(X\) is projective, GAGA identifies this space with the usual Berkovich
analytification over trivially valued \(\mathbf C\)
\cite[Remark~1.3.3]{MP24}.

For a coherent ideal \(\mathfrak a\), set

\[
              \log|\mathfrak a|(v):=-v(\mathfrak a)
              \in[-\infty,0].
\tag{59.3}\label{eq:mp-log-ideal}
\]

Thus

\[
 \log|\mathfrak a\mathfrak b|
   =\log|\mathfrak a|+\log|\mathfrak b|,
 \qquad
 \log|\mathfrak a+\mathfrak b|
   =\max\{\log|\mathfrak a|,\log|\mathfrak b|\}.
\tag{59.4}\label{eq:mp-tropical-ideal-laws}
\]

The value \(-\infty\) is essential: it records that \(\mathfrak a\) vanishes
identically on the support of a semivaluation.  Real-valued model functions
are obtained by truncating before taking rational linear combinations.

\begin{definition}[Piecewise-linear functions]
\label{def:mp-pl-functions}
Let \(\PL^+(X^{\bethan})\) consist of functions
\[
  \phi=\frac1m\max_{1\leq j\leq N}
       \{\log|\mathfrak a_j|+k_j\},
  \qquad m\geq1,\quad k_j\in\mathbf Z,
\tag{59.5}\label{eq:mp-pl-plus}
\]
whose displayed maximum is finite everywhere.  Let
\(\PL(X^{\bethan})\) be the \(\mathbf Q\)-vector space generated by
\(\PL^+(X^{\bethan})\), and put
\(\PL_{\mathbf R}:=\PL\otimes_{\mathbf Q}\mathbf R\).
\end{definition}

\begin{proposition}\label{prop:mp-pl-dense}
The space \(\PL(X^{\bethan})\) is uniformly dense in
\(C^0(X^{\bethan})\).
\end{proposition}

\begin{proof}
If \(v\ne w\), choose \(\mathfrak a\) with
\(v(\mathfrak a)\ne w(\mathfrak a)\), and choose a rational number strictly
between these two values.  A truncation of \(\log|\mathfrak a|\) then
separates \(v\) and \(w\).  The space \(\PL\) contains the rational
constants and is stable under maxima and minima.  The lattice form of the
Stone--Weierstrass theorem proves the claim.  This is
\cite[Lemma~1.4.3 and Proposition~1.4.4]{MP24}.
\end{proof}

For a coherent ideal $\mathfrak a$, its integral closure
$\overline{\mathfrak a}$ is defined stalkwise by monic equations
\[
 f^d+c_1f^{d-1}+\cdots+c_d=0,
 \qquad c_i\in\mathfrak a^i,
\]
as in \cref{def:appA-integral-closure-ideal}. It is a coherent ideal
by \cite[Theorem~2.9]{LejeuneTeissier08}. For this analytic input,
a nonzero ideal on our irreducible space defines a nowhere dense
subspace. The zero and unit ideals have integral closures $0$ and
$\mathcal O_X$, respectively.

For $\mathfrak a\ne0$, let $\nu:Y\to X$ be its normalized blowup.
Then $Y$ is normal and the extended ideal defines an effective
Cartier divisor $E$ with
\[
 \mathfrak a\cO_Y=\cO_Y(-E),\qquad
 \nu_*\cO_Y(-E)=\overline{\mathfrak a}.
\tag{59.6}\label{eq:mp-normalized-blowup}
\]
The analytic integral-dependence theorem gives
$\overline{\mathfrak a}=\nu_*(\mathfrak a\cO_Y)\cap\cO_X$
\cite[Theorem~2.1 and the proof of Theorem~2.9]{LejeuneTeissier08}.
Normality of $X$ gives $\nu_*\cO_Y=\cO_X$, removing the intersection
and proving \eqref{eq:mp-normalized-blowup}. The same argument
applies to a normal proper modification on which the ideal becomes
invertible. The divisor $E$ need not be exceptional: for an
invertible ideal $\cO_X(-D)$ the blowup is the identity and $E=D$.
The zero ideal is excluded from this blowup construction.

\begin{proposition}\label{prop:mp-integral-closure-values}
For every coherent ideal $\mathfrak a$ and every
$v\in X^{\bethan}$,
\[
                     v(\mathfrak a)=v(\overline{\mathfrak a}).
\tag{59.7}\label{eq:mp-integral-closure-values}
\]
For coherent ideals $\mathfrak a,\mathfrak b$,
\[
 \overline{\mathfrak a}\subseteq\overline{\mathfrak b}
 \quad\Longleftrightarrow\quad
 v(\mathfrak a)\geq v(\mathfrak b)
 \quad\text{for every }v\in X^{\divis}.
\]
In particular, equality of the ideal functions on the divisorial
locus determines the integral closure. Here $X^{\divis}$ includes
$v_{\triv}$, as in \cref{def:mp-divisorial-locus}.
\end{proposition}

\begin{proof}
The full-support case of \eqref{eq:mp-integral-closure-values} is
\cite[Lemma~2.2.4]{MPTropicalV2}. To include infinite values, put
$\mathfrak j=\overline{\mathfrak a}$. We may assume
$\mathfrak a\ne0$. Locally, finite generation and the displayed
integral equations make $\mathcal O_X[\mathfrak jT]$ a finite
graded module over $\mathcal O_X[\mathfrak aT]$; see
\cref{lem:appA-rees-translation} and \cite[Tag~00H3]{stacks-project}.
In degrees larger than those of homogeneous module generators,
this gives $\mathfrak j^{N+1}=\mathfrak a\mathfrak j^N$.
The equality persists for all larger exponents. A finite cover of
the compact space $X$ therefore gives a single $N\geq1$ such that
\[
           \mathfrak j^{N+1}=\mathfrak a\mathfrak j^N
           \quad\text{on }X.
\]
Apply the character $v$ only to these globally defined ideals.
If $v(\mathfrak j)=+\infty$, the inclusion
$\mathfrak a\subseteq\mathfrak j$ forces
$v(\mathfrak a)=+\infty$. Otherwise the identity gives
\[
 (N+1)v(\mathfrak j)=v(\mathfrak a)+Nv(\mathfrak j).
\]
Its left side is finite, so cancellation is legitimate and proves
\eqref{eq:mp-integral-closure-values}. This argument does not
require an arbitrary character to act on local holomorphic germs.

For the converse criterion, first assume $\mathfrak b\ne0$ and
take its normalized blowup $\nu:Y\to X$. Write
\[
             \mathfrak b\cO_Y=\cO_Y(-D),
             \qquad D=\sum_{i=1}^r b_iE_i,
             \qquad b_i>0.
\]
There are finitely many components since $Y$ is compact. Their
generic orders on pulled-back coherent ideals define divisorial
valuations on $X$. The following finite test suffices:
\[
 \mathfrak a\subseteq\overline{\mathfrak b}
 \quad\Longleftrightarrow\quad
 \ord_{E_i}(\mathfrak a)\geq b_i
 \quad(1\leq i\leq r).
\]
Indeed, if these inequalities hold, divide a local generator of
$\mathfrak a\cO_Y$ by a local equation of $D$.
The generic order of a coherent ideal is constant off a proper
analytic subset of each $E_i$, so its order along every local
branch is the same. The quotient therefore has nonnegative order
along each branch; away from $D$ the denominator is a unit. A meromorphic function on
a normal space with no codimension-one pole is holomorphic.
This follows stalkwise from \cite[Tag~031T]{stacks-project};
the analytic divisor argument is also given in
\cite[Proposition~4.8]{LejeuneTeissier08}.
Thus $\mathfrak a\cO_Y\subseteq\cO_Y(-D)$, and pushing forward
gives $\mathfrak a\subseteq\overline{\mathfrak b}$ by
\eqref{eq:mp-normalized-blowup}. The reverse implication follows
from integral-closure invariance and ideal order reversal. These
steps give the coherent-ideal form of the local analytic criterion
in \cite[Theorems~2.1 and~4.6]{LejeuneTeissier08}.

If $v(\mathfrak a)\geq v(\mathfrak b)$ for every divisorial $v$,
the finite test gives $\mathfrak a\subseteq\overline{\mathfrak b}$,
hence $\overline{\mathfrak a}\subseteq\overline{\mathfrak b}$ by
idempotence of integral closure. Conversely, this inclusion and
\eqref{eq:mp-integral-closure-values} give all the claimed
inequalities. If $\mathfrak b=\cO_X$, then $D=0$ and the finite
test is empty, as it should be. If $\mathfrak b=0$ and
$\mathfrak a\ne0$, the trivial valuation has
$v_{\triv}(\mathfrak a)=0<v_{\triv}(0)=+\infty$, so the criterion
fails on both sides. If $\mathfrak a=0$, both sides hold for any
$\mathfrak b$. These cases also cover a zero-dimensional $X$.
\end{proof}

\begin{example}\label{ex:mp-integral-closure-divisor-tests}
Let $p=[0:0:1]\in\mathbf P^2$, with affine coordinates $x,y$ near
$p$. Let $\mathfrak a$ be the ideal generated by the homogeneous
sections $X_0^2,X_1^2$ of $\cO_{\mathbf P^2}(2)$, viewed as the
image of $\cO(-2)^{\oplus2}\to\cO$, and put
$\mathfrak m=\mathcal I_p$. Then
$\mathfrak a_p=(x^2,y^2)$ and
$(\mathfrak m^2)_p=(x^2,xy,y^2)$; both ideals are the unit ideal
away from $p$.

Blow up $p$. On the chart $x=u$, $y=ut$, the pullbacks of these
two ideals are both $(u^2)$; the other chart gives the same
conclusion. Thus both pullbacks are $\cO(-2E)$, where $E$ is
the exceptional divisor. A germ pulls back to a section of
$\cO(-2E)$ exactly when its Taylor series has order at least two.
Consequently
\[
        \overline{\mathfrak a}=\mathfrak m^2\ne\mathfrak a,
        \qquad
        \ord_E(\mathfrak a)=\ord_E(\mathfrak m^2)=2.
\]
The strict inclusion is also seen from $xy\notin(x^2,y^2)$,
while $(xy)^2-x^2y^2=0$ is an integral equation.
The exact reduction identity is
$(\mathfrak m^2)^2=\mathfrak a\mathfrak m^2$.
The character supported at $p$ takes value $+\infty$ on both
ideals, illustrating the infinite-value case of the proposition.

Every prime divisor on $\mathbf P^2$ itself gives value zero to
both $\mathfrak m$ and $\mathfrak m^2$, whereas $\ord_E$ gives
values one and two. Thus divisors on $X$ alone do not suffice.
Conversely, exceptional divisors alone do not suffice either:
on $\mathbf P^1$ the ideals $\cO(-p)$ and $\cO(-2p)$ are distinct
and integrally closed, their blowups are the identity, and their
different values are detected by the nonexceptional divisor $p$.
\end{example}


\begin{warning}{Valuations versus semivaluations}
A bimeromorphic modification \(\pi:Y\to X\) induces a bijection
\(Y^{\val}\simeq X^{\val}\) on full-support valuations
\cite[Proposition~2.1.6]{MP24}.  It need not induce a bijection on all of
\(Y^{\bethan}\): proper-support semivaluations remember subspaces that may be
created or contracted by \(\pi\).
\end{warning}

\section{The Stein-local algebraic interface}

We isolate the analytic--algebraic passage used below.  If \(S\) is a complex
space and \(p\in S\), let \(\operatorname{ProjGerm}(S,p)\) be the category
whose objects are projective morphisms \(Z_U\to U\), for neighborhoods
\(p\in U\subset S\), and whose morphisms are identified after restriction to
a smaller neighborhood.  Thus this category remembers only the germ over
\((S,p)\).

\begin{theorem}\label{thm:mp-bingener-projective-germs}
Let \(S\) be a complex space and \(p\in S\).  There is an equivalence
\[
 \left\{
 \begin{array}{c}
   \text{projective schemes of finite presentation}\cr
   \text{over }\operatorname{Spec}\cO_{S,p}
 \end{array}
 \right\}
 \simeq \operatorname{ProjGerm}(S,p).
\tag{59.7a}\label{eq:mp-projective-germ-equivalence}
\]
It is compatible with closed immersions, blowups along ideals of finite
type, and composition of projective morphisms.
\end{theorem}

\begin{proof}
Choose a Stein neighborhood \(U\) of \(p\), put
\(A_U=\Gamma(U,\cO_U)\), and regard the singleton \(K=\{p\}\) as a
Stein compactum.  It is semianalytic and hence excellent by
\cite[Corollary~4.4.5]{XiaStein26}; moreover
\[
              \Gamma(K,\cO_U)=\cO_{S,p}
              =\varinjlim_{p\in U'}\Gamma(U',\cO_{U'}),
\]
where \(U'\) runs over Stein neighborhoods of \(p\).  A projective
\(\cO_{S,p}\)-scheme is of finite presentation, since \(\cO_{S,p}\) is
noetherian, and admits a closed immersion into
\(\mathbf P^N_{\cO_{S,p}}\).  Its homogeneous ideal and the finitely many
relations needed to define the quotient descend to \(A_{U'}\) after
shrinking.  Bingener's relative analytification
\cite[Theorem~5.3.1]{XiaStein26} turns the descended closed subscheme into
a closed analytic subspace of \(U'\times\mathbf P^N\), hence into an object
of \(\operatorname{ProjGerm}(S,p)\).

Conversely, after shrinking, an analytic projective germ admits a closed
immersion into \(U\times\mathbf P^N\).  Apply Bingener's algebraization
equivalence to its coherent ideal sheaf on the analytification of
\(\mathbf P^N_{A_U}\).  The ambient projective space is proper over the
base, so every coherent module has proper support and
\cite[Theorem~5.4.4 and Corollary~5.4.12]{XiaStein26} apply with
\(K=\{p\}\).  The algebraized quotient of
\(\cO_{\mathbf P^N_{\cO_{S,p}}}\) defines the required projective scheme.
Corollary~5.4.13 of the same source shows directly that closed subspaces
correspond, including their quotient-algebra structures.

For full faithfulness, a morphism between projective objects is equivalent
to its graph in their fiber product.  The graph has proper support, and
\cite[Corollary~5.4.14]{XiaStein26} identifies algebraic morphisms over
\(\cO_{S,p}\) with analytic morphism germs.  Hence the two object
constructions are quasi-inverse functors.  Closed immersions and composition
are preserved by construction.  A blowup is the relative Proj of the Rees
algebra of a finite-type ideal, and relative analytification commutes with
finite products and closed subspaces; it therefore preserves these blowups.

Relative analytification and its flat, completion-preserving local-ring map
are proved in \cite[Section~1]{Bingener76Stein} and reconstructed in
\cite[Theorem~5.3.1]{XiaStein26}.  The coherent algebraization used in the
reverse direction is Bingener's Aussagen~5.1 and~5.3; see
\cite[Theorem~5.4.4 and Corollaries~5.4.12--5.4.14]{XiaStein26}.
\end{proof}

The dimension assertion needed in the birational argument is more delicate
than the mere existence of the equivalence.

\begin{proposition}\label{prop:mp-bingener-dimension}
Let \(U\) be Stein, let \(p\in U\), and let \(Z\to U\) be projective.
Denote by \(Z_p^{\mathrm{alg}}\) the projective
\(\cO_{U,p}\)-scheme corresponding to its germ under
\Cref{thm:mp-bingener-projective-germs}.  After replacing \(U\) by a
sufficiently small Stein neighborhood of \(p\),
\[
                 \dim Z_p^{\mathrm{alg}}=\dim Z|_U.
\tag{59.7b}\label{eq:mp-bingener-dimension}
\]
More precisely, after shrinking there are closed constructible subsets
\(T_m\) of the algebraic model such that \(T_m^{\mathrm{an}}\) is the locus where
the analytic local dimension is at least \(m\), while
\(T_m\otimes_A\cO_{U,p}\) is the union of the irreducible components of
\(Z_p^{\mathrm{alg}}\) of dimension at least \(m\).
\end{proposition}

\begin{proof}
Choose an algebraic projective model over the Stein algebra
\(A=\Gamma(U,\cO_U)\), as in
\Cref{thm:mp-bingener-projective-germs}.  This is Bingener's
Aussage~2.8 with the Stein compactum \(K=\{p\}\)
\cite[Aussage~2.8]{Bingener76Stein}.  The proof constructs closed
constructible sets \(T_n\) whose analytifications are the loci where the
analytic local dimension is at least \(n\), while
\(T_n\otimes_A A_K\), with \(A_K=\cO_{U,p}\), is the union of the
algebraic components of dimension at least \(n\).  Taking the largest
\(n\) for which these sets are nonempty gives
\eqref{eq:mp-bingener-dimension}.  Properness and finite presentation are
both essential.  Bingener's counterexample immediately after Aussage~2.8
is \(S=\mathbf C\), \(X=\operatorname{Spec}A_z\), and \(K=\{0\}\): it has
algebraic germ dimension zero but analytic dimension one on every
neighborhood of \(0\).
\end{proof}

The point of \Cref{prop:mp-bingener-dimension} is the following.  A
divisorial valuation can be made into a divisor after finitely many blowups
of its centers over the local scheme \(\operatorname{Spec}\cO_{S,p}\).
The proposition then shows that the corresponding irreducible analytic
center has codimension one after shrinking.  Without this last step, a
divisor over the local scheme would not yet be known to come from a divisor
on an analytic modification.

Here is the simplest local test of the interface.  Suppose that \(S\) is
smooth at \(p\), with parameters \(z_1,\ldots,z_d\), and let
\(I=(z_1,\ldots,z_r)\subset\cO_{S,p}\).  The algebraic blowup is covered by
the usual affine charts
\[
 \operatorname{Spec}\cO_{S,p}
   [z_1/z_j,\ldots,z_r/z_j],
 \qquad 1\leq j\leq r.
\]
After shrinking to a Stein coordinate neighborhood, relative
analytification gives the corresponding analytic charts of
\(\operatorname{Bl}_{V(I)}S\).  The exceptional fiber is
\(\mathbf P^{r-1}\), both algebraically and analytically, and the ambient
dimension remains \(d\).  This calculation illustrates the two logically
separate inputs above: compatibility with the Rees construction identifies
the blowup, while Aussage~2.8 controls its dimension after shrinking.

Nothing in this comparison creates positivity.  The hypotheses
``finite presentation'', ``projective'' (or coherent sheaf with proper
support), and ``excellent Stein compactum'' are precisely what make the
algebraization step available.  The later nef, Monge--Amp\`ere, and envelope
arguments come from \cite{MP24,MPWN25}, not from Bingener's theorem.

\section{Divisorial and monomial valuations}

Let \(\pi:Y\to X\) be a proper bimeromorphic morphism with \(Y\) normal,
and let \(F\subset Y\) be a prime divisor.  For a coherent ideal
\(\mathfrak a\) on \(X\), define

\[
 \ord_F(\mathfrak a)
 :=\ord_F(\mathfrak a\cO_Y).
\tag{59.8}\label{eq:mp-divisorial-order}
\]

This is independent of the generic point used to compute it.  The
birational invariance just noted lets us regard it as a valuation on \(X\).

\begin{definition}[Divisorial locus]
\label{def:mp-divisorial-locus}
Following the rational convention of \cite[Definition~2.3.1]{MP24}, put
\[
 X^{\divis}:=\{r\ord_F\mid r\in\mathbf Q_{>0},\ F
       \text{ a prime divisor over }X\}\cup\{v_{\triv}\}.
\tag{59.9}\label{eq:mp-divisorial-locus}
\]
Here \(v_{\triv}\) is zero on every nonzero coherent ideal.
\end{definition}

For monomial valuations, take a smooth proper bimeromorphic model
$\pi:Y\to X$ and a reduced snc divisor $B=\sum_{i\in I}B_i$ on $Y$.
For $J\subseteq I$, put $B_J=\bigcap_{i\in J}B_i$, with
$B_\varnothing=Y$. Each connected component $Z$ of $B_J$ gives a cone
\[
                         \widehat\sigma_Z=(\mathbf R_{\geq0})^J.
\tag{59.10}\label{eq:mp-monomial-cone}
\]
If $K\subseteq J$, identify the face on which $w_i=0$ for
$i\in J\setminus K$ with the cone of the unique connected component
of $B_K$ containing $Z$. These identifications define
$\widehat\Delta(Y,B)$. Different components of the same intersection
give different cones, even though their weight coordinates have the
same indices. Compactness of $Y$ makes this a finite cone complex.

Choose $p\in Z$ and coordinates centered at $p$ with
$B_i=(z_i=0)$ for $i\in J$. Give every other coordinate weight zero.
For $f\in\cO_{Y,p}$, define
\[
 \val_{w,p}\!\left(\sum_{\alpha\in\mathbf N^n}c_\alpha z^\alpha\right)
   :=\min_{c_\alpha\ne0}\sum_{i\in J}\alpha_iw_i,
 \qquad \val_{w,p}(0)=+\infty.
\tag{59.11}\label{eq:mp-monomial-local}
\]
For a coherent ideal $\mathfrak a$ on $X$, evaluate the finitely
generated ideal $(\mathfrak a\cO_Y)_p$ by taking the minimum of
the values of local generators.

\begin{proposition}\label{prop:mp-monomial-globalization}
This ideal value is independent of $p\in Z$ and of the adapted
coordinates. It defines a full-support valuation on $X$, and the
face identifications give a continuous injection
\[
 \widehat\Delta(Y,B)\hookrightarrow X^{\val},
 \qquad w\longmapsto\val(w).
\tag{59.12}\label{eq:mp-cone-complex-embedding}
\]
If $K=\{i\in J:w_i>0\}$, the corresponding valuation on $Y$ has
center the component $W$ of $B_K$ containing $Z$. For $K=\varnothing$
it is $v_{\triv}$. Every rational point of the cone complex is
divisorial in the convention of \cref{def:mp-divisorial-locus}.
\end{proposition}

\begin{proof}
The analytic globalization is
\cite[Proposition~2.3.4]{MPTropicalV2}; the algebraic coordinate
construction is \cite[Proposition~3.1]{JM12}. We give the analytic
argument, including the faces and the different intersection components.

First separate the coordinates of positive weight, denoted $u$, from
the remaining coordinates $\zeta$, and write
\[
                  f(u,\zeta)=\sum_\beta f_\beta(\zeta)u^\beta.
\]
For a nonzero germ, its value is the minimum of
$\sum_{i\in K}w_i\beta_i$ for which $f_\beta$ is a nonzero germ.
There are only finitely many positive-weight exponents below any
fixed bound, so the minimum is attained. The initial form of $f$
is a nonzero polynomial in $u$ over the domain $\mathbf C\{\zeta\}$.
Initial forms multiply without becoming zero. This proves
multiplicativity; the ultrametric inequality follows from addition
of the expansions. If there is no positive weight, every nonzero
germ has value zero. In all cases local units have value zero, and
the minimum on an ideal is independent of its finite generators.

For coordinate independence, fix $s\geq0$ and consider the ideal
\[
 \mathfrak f_w^s
   =\sum_{\substack{\alpha\in\mathbf N^J\\
                     \langle w,\alpha\rangle\geq s}}
           \cO_Y\!\left(-\sum_{i\in J}\alpha_i B_i\right).
\]
The upward-closed set of exponents has finitely many minimal
elements: its monomial ideal in a polynomial ring is finitely
generated. Hence $\mathfrak f_w^s$ is coherent. In adapted
coordinates its stalk is precisely the monomial ideal generated by
these minimal exponents, and
\[
          \val_{w,p}(f)\geq s
          \quad\Longleftrightarrow\quad
          f\in(\mathfrak f_w^s)_p.
\]
Indeed, membership means that every nonzero monomial of the power
series is divisible by one of the generators. Grouping terms
according to a chosen divisor monomial gives convergent quotients
on a smaller polydisc. Since $\mathfrak f_w^s$ uses only the
divisor ideals, this criterion proves coordinate independence,
also when some weights vanish.

To move $p$ along $Z$, use a product polydisc with normal
coordinates $z'=(z_i)_{i\in J}$ and tangential coordinates $z''$.
For a holomorphic function on this polydisc write
\[
                     f=\sum_\alpha g_\alpha(z'')(z')^\alpha.
\]
A coefficient function $g_\alpha$ is either identically zero or
has a nonzero germ at every point of the connected tangential
polydisc, by the identity theorem. Its numerical value may vanish
at a point without its germ vanishing. Thus the set of exponents
contributing to \eqref{eq:mp-monomial-local} is unchanged as
$p$ moves. Coherence supplies generators on a common neighbourhood,
so the same conclusion holds for each coherent ideal. Connectedness
of $Z$ gives global independence. Nonzero coherent ideals on the
connected manifold $Y$ have nonzero stalks and therefore finite
values. We obtain a character in $Y^{\val}$ and then one in
$X^{\val}$ by pullback of ideals.

Omitting the zero-weight normal coordinates leaves the same formula.
Applying point independence along the larger stratum $W$ proves
compatibility with every face identification. If $K\ne\varnothing$,
an ideal has positive value exactly when its image in $\cO_W$
is zero: this is the vanishing of the coefficient of $u^0$ in all
its local generators. This identifies the center as $W$.

To prove injectivity, first work on $Y$ and extend the weight vector
by zero to $I$. Point independence lets us compute at a point of
$W$ outside the $B_i$ with $i\notin K$; such points exist by the
snc condition. The global ideal $\cO_Y(-B_i)$ then has value $w_i$
for every $i\in I$. Thus distinct extended vectors are separated.
Suppose that the extended vectors agree, with nonempty positive
support $K$, but their centers are different components $W,W'$ of
$B_K$. These components are disjoint. Their global ideal sheaves
give
\[
     \val_{W,w}(\mathcal I_W)=\min_{i\in K}w_i>0,
     \qquad \val_{W',w}(\mathcal I_W)=0.
\]
Hence these valuations are different as well. If both the vector
and its component agree, the face identifications already identify
the two cone points. The zero vectors all give the common vertex.
Finally, \cref{prop:ts-bimeromorphic-valuations} transports
injectivity from $Y^{\val}$ to $X^{\val}$.

For continuity on a fixed cone, choose one point of its stratum and
generators of a fixed nonzero ideal there. The union of their power
series supports has finitely many coordinatewise minimal exponents,
again by finite generation of monomial ideals. Since all weights
are nonnegative, the ideal value is the minimum of the corresponding
finitely many linear forms in $w$. It is continuous on the whole
closed cone, including its faces. The zero-ideal coordinate is
constant $+\infty$. The quotient topology of the cone complex and
pointwise convergence on ideals prove continuity of
\eqref{eq:mp-cone-complex-embedding}. Compare the algebraic cone
description in \cite[Lemma~4.5 and the following paragraph]{JM12}.

For a nonzero rational vector, omit its zero coordinates and write
$w=ca$, where $c\in\mathbf Q_{>0}$ and
$a=(a_1,\ldots,a_k)\in\mathbf Z_{>0}^k$ is primitive.
Here $k=|K|$. If $k=1$, then $a_1=1$ and the valuation is
$c\ord_W$. If $k\geq2$, blow up the smooth stratum $W$.
Choose $j$ with $a_j=\min_i a_i$. In the chart with coordinates
$u_j$ and $u_i/u_j$ for $i\ne j$, the weights become
\[
                    a_j,\qquad a_i-a_j\quad(i\ne j).
\]
Substitution preserves every monomial's weight: it replaces
the exponent of $u_j$ by the sum of the normal exponents and
leaves the other exponents unchanged. This exponent map is
injective, so the resulting monomial valuation pulls back the
original ideal character. The same charts show that the reduced
total transform of $B$ is snc. Discard the zero weights and take
the component containing the local positive-weight stratum as
the next center. The construction is again on a global smooth model.

The sum of the positive integer weights decreases by
$(k-1)a_j>0$, while their greatest common divisor stays one.
Repeating the construction terminates with one weight, equal to
one. On that model the valuation is $\ord_F$ for a prime divisor
$F$, and the original rational valuation is $c\ord_F$.
The zero vector gives $v_{\triv}$. These calculations give the
rational realization asserted in
\cite[Remark~2.3.5]{MPTropicalV2}, with its scale $c$ retained.
\end{proof}

\begin{example}\label{ex:mp-monomial-zero-weights}
Near a smooth divisor $(u=0)$, give $(u,t)$ weights $(a,0)$ with
$a>0$. The germ $ut+u^2$ has value $a$ at every point of the
divisor, including $t=0$: the coefficient germ $t$ is nonzero
there. Substituting the number $t=0$ before taking the weighted
order would instead give $2a$. Zero-weight coordinates are kept
as coefficient variables, not evaluated at the chosen point.

At a smooth point of dimension $n\geq2$, the vector
$(1,\ldots,1)$ gives the order along the exceptional divisor of
the ordinary blowup, equivalently the maximal-ideal order. In
dimension one the blowup of the point ideal is the identity, and
the same order is $\ord_p$ on the curve itself. These constructions
use coherent ideals even when $X$ has no nonconstant meromorphic
functions.
\end{example}

\begin{example}\label{ex:mp-monomial-distinct-strata}
On $Y=\mathbf P^1\times\mathbf P^1$, with affine coordinates
$(x,y)$, take $B_1=(y=0)$ and let $B_2$ be the graph of the
morphism $x\mapsto x^2-1$. Both curves are smooth and they meet
transversely at $p_+=(1,0)$ and $p_-=(-1,0)$. The cone complex
has two quadrants, whose corresponding boundary rays are
identified. Their interior points of weight $(1,1)$ give the
orders of the exceptional divisors over $p_+$ and $p_-$.
Both assign value one to the ideals of $B_1$ and $B_2$, but they
assign values one and zero, respectively, to $\mathcal I_{p_+}$.
This is why the intersection component is part of the cone datum.
On the face $(a,0)$ the two quadrants both give $a\ord_{B_1}$.

For the weight $(2/3,1)=\frac13(2,3)$ at $p_+$, successive
blowups transform the integer weights as
\[
        (2,3)\longmapsto(2,1)\longmapsto(1,1)
             \longmapsto(1,0).
\]
With $u=y$ and $v=y-(x^2-1)$, the resulting prime divisor has
$\ord_F(u)=2$ and $\ord_F(v)=3$. Thus
$\val(2/3,1)=\frac13\ord_F$, and the local ideal $(u^3,v^2)$
has value two. The value $\val(2/3,1)(u)=2/3$ rules out an
unscaled prime order.
\end{example}

\section{Models, dual complexes, and retractions}

A \emph{model} of $X$ is a normal compact K\"ahler space $\cX$, a
flat $\mathbf C^*$-equivariant morphism $\cX\to\mathbf P^1$, and a
$\mathbf C^*$-equivariant bimeromorphic morphism
\[
       \mu:\cX\longrightarrow X\times\mathbf P^1
\tag{59.13}\label{eq:mp-analytic-model}
\]
that is a biholomorphism over $\mathbf P^1\setminus\{0\}$.
Here $\mathbf C^*$ acts trivially on $X$ and by multiplication
on the affine coordinate $t$ of $\mathbf P^1$. A model is snc
when $\cX$ is smooth and $\cX_{0,\mathrm{red}}$ is an snc divisor.
We use the cofinal directed system of projective snc models when
applying local algebraization. Equivariant resolution gives this
cofinality; see \cite[Remark~2.5.2]{MPTropicalV2}.

Write
\[
                         \cX_0=\sum_{i\in I}b_iE_i,
                         \qquad b_i=\ord_{E_i}(t)>0.
\tag{59.14}\label{eq:mp-central-fiber}
\]
The connected group $\mathbf C^*$ preserves each of the finitely
many irreducible components. Each component therefore defines an
invariant normalized vertical valuation
\[
 v_{E_i}:=b_i^{-1}\ord_{E_i},\qquad v_{E_i}(t)=1.
\tag{59.14a}\label{eq:mp-normalized-vertical-valuation}
\]
Let $p_1:X\times\mathbf P^1\to X$ be the projection. Pullback of
ideals under $p_1\circ\mu$ defines
\[
 r(v_{E_i})(\mathfrak a)
   :=b_i^{-1}\ord_{E_i}
      \bigl((p_1\circ\mu)^{-1}\mathfrak a\cdot\cO_{\cX}\bigr),
 \qquad \mathfrak a\in\mathfrak I_X.
\tag{59.14b}\label{eq:mp-vertical-restriction}
\]
Thus restriction means restriction to functions pulled back from
$X$, not passage to the quotient $t=0$. A valuation factoring
through that quotient would assign $+\infty$ to $t$, whereas
$v_{E_i}(t)=1$. Denote the set of these normalized valuations,
over all models, by $\mathcal D_{\mathbf C^*}(X)$.

The local algebra in the comparison requires the relative
dimension formula. Discreteness and the normalization at $t$
alone do not establish the required geometric realization.

\begin{lemma}\label{lem:mp-local-vertical-restriction}
Let $A$ be a regular excellent local $\mathbf C$-algebra,
$K=\operatorname{Frac}A$, and let $w$ be a positive rational
multiple of a prime-divisor order over $\operatorname{Spec}A[t]$.
Assume that $w$ is centered on $\operatorname{Spec}A[t]$ and
$w(t)>0$. Then $v=w|_K$ is either trivial or a positive rational
multiple of a prime-divisor order over $\operatorname{Spec}A$.
\end{lemma}

\begin{proof}
Assume that $v$ is nontrivial. Write $\mathfrak p$ for its center
on $A$ and set $h=\dim A_{\mathfrak p}$. The center of $w$ on
$A[t]$ is $\mathfrak q=(\mathfrak p,t)$, of height $h+1$.
The dimension formula for a divisor on a birational model gives
$\operatorname{trdeg}_{\kappa(\mathfrak p)}\kappa(w)=h$.

Here $\kappa(w)$ and $\kappa(v)$ are the residue fields of the
valuation rings. The field extension $K(t)/K$ has transcendence
degree one, so
$\operatorname{trdeg}_{\kappa(v)}\kappa(w)\leq1$.
Indeed, lifts of residue elements algebraically independent over
$\kappa(v)$ are algebraically independent over $K$: divide any
putative polynomial relation by a coefficient of least $v$-value
and reduce in $\kappa(w)$. At least one reduced coefficient is
nonzero, giving the contradiction.

Since $\Gamma_v$ is a nonzero subgroup of the discrete cyclic
group $\Gamma_w$, its rational rank is one. Abhyankar's
inequality on the regular excellent scheme $\operatorname{Spec}A$
therefore gives
\[
 h-1\leq\operatorname{trdeg}_{\kappa(\mathfrak p)}\kappa(v)
       \leq h-1.
\]
Thus $v$ is Abhyankar of rational rank one, hence divisorial by
\cite[Proposition~3.7 and Remark~3.9]{JM12}. Its scale is rational
because $\Gamma_v\subseteq\Gamma_w\subset\mathbf Q$.
The dimension formula and inequality used here are recalled in
\cite[Section~3.2, equation~(3.4)]{JM12}. This is the relative
local-ring form of \cite[Lemma~4.1]{BHJ17}; it does not assume
that $K$ is finitely generated over $\mathbf C$.
\end{proof}

\begin{theorem}\label{thm:mp-invariant-divisorial-comparison}
The map defined by \eqref{eq:mp-vertical-restriction} satisfies
\[
                       r\bigl(\mathcal D_{\mathbf C^*}(X)\bigr)
                          =X^{\divis}.
\tag{59.14c}\label{eq:mp-invariant-divisorial-comparison}
\]
The trivial valuation is obtained from
$X\times\{0\}\subset X\times\mathbf P^1$.
\end{theorem}

\begin{proof}
This is \cite[Theorem~2.5.5]{MPTropicalV2}; the algebraic
comparison is \cite[Theorem~4.6]{BHJ17}. We give the local
restriction and the global extraction separately.

\emph{Local restriction.}
Take a projective snc model and a component $E$ of its central
fiber. Put $b=\ord_E(t)$ and $v=r(b^{-1}\ord_E)$.
The morphism $\cX\to X$ is projective, since both $\mu$ and
$X\times\mathbf P^1\to X$ are projective. At a point $p$ of
the image of $E$, \cref{thm:mp-bingener-projective-germs}
therefore algebraizes this morphism, its map to
$\mathbf P^1_A$, and its divisor, where $A=\cO_{X,p}$.
The local rings of $X$ are regular and excellent; excellence
for convergent power series is recalled in \cite[Section~1.1]{JM12}.
Choose a prime component of the algebraized divisor. The
completion-preserving local comparison preserves normality and
its divisor order. The algebraized morphism is birational to
$\mathbf P^1_A$, since it is an isomorphism away from $t=0$.
The component's order divided by $b$ therefore gives a
divisorial valuation $w$ on $K(t)$, where
$K=\operatorname{Frac}A$, centered over $t=0$.

The zero ideal has value $+\infty$ in both descriptions. For a
nonzero coherent ideal, its order can be computed on the chosen
local component. Factor the generic divisor order from the
pulled-back ideal along the smooth divisor $E$. The remaining
ideal has a nonzero restriction to $E$. It cannot vanish
identically on any nonempty analytic open subset of the
irreducible $E$. The same order is therefore read on any
component occurring in the germ. This is the order comparison
in \cite[Proposition~2.5.10]{MPTropicalV2}, with local branches
retained. We do not assume that a globally irreducible center
has an irreducible germ.

By \cref{lem:mp-local-vertical-restriction}, $w|_K$ is trivial
or rational divisorial. If it is trivial, the order comparison
gives $v=v_{\triv}$. Otherwise its valuation ring is a DVR
essentially of finite type over $A$.

\emph{Extraction by global centers.}
We recall why successive center blowups terminate for a
divisorial valuation. Normalize its value group to $\mathbf Z$.
At each stage let $R$ be the local ring at its center and choose
$u$ among a finite set of generators of its maximal ideal,
with $u$ having least positive value.
The blowup chart $R[\mathfrak m_R/u]$, localized at the new
center, is contained in the valuation ring. If $a/b$ has
nonnegative value and $b$ has positive value, then $a,b\in
\mathfrak m_R$ and this chart contains $a/u$ and $b/u$.
The value of the denominator has decreased by a positive integer.
Repeating shows that every fixed fraction in the valuation ring
belongs to a center local ring after finitely many steps.

The divisorial DVR is a localization of a finitely generated
algebra over the initial ring. Apply the preceding argument to
its finitely many algebra generators. Elements inverted in that
localization have value zero, so are already units in the center
local ring. At a finite stage the center local ring is therefore
the DVR itself. Normalizing the intermediate blowups does not
change this argument: the local rings remain contained in the
same DVR. This proves the center-blowup assertion used in
\cite[Lemma~2.5.9]{MPTropicalV2} and
\cite[proof of Lemma~4.7]{BHJ17}.

Apply this procedure on $X$ to the global center of $v$, and
then to its centers on the successive normalized blowups.
These are irreducible closed analytic subspaces: on a common
resolution they are proper images of the transform of $E$.
Their reduced ideal sheaves are coherent. All blowups are
projective and are made globally. Localize the construction at
the fixed point $p$. Compatibility with blowups identifies the
local center rings with the preceding algebraic process.
Proper images commute with this germ comparison, so the chosen
local divisor component dominates a component of each center
germ. If a global center splits into several germ components, its reduced
ideal becomes the ideal of the selected component after
localizing at that component's generic point; the other
components do not affect this center local ring.

The algebraic process terminates with a codimension-one center.
By \cref{prop:mp-bingener-dimension}, the corresponding analytic
center has dimension $\dim X-1$ near $p$. Its global
irreducibility makes it a prime divisor $F$ on the resulting
normal modification. The local DVR equality and the order
comparison above give $v=q\ord_F$ on every coherent ideal.
Here $q=e/b$ for some positive integer $e$, since
$\Gamma_{w|_K}$ is a nonzero subgroup of $b^{-1}\mathbf Z$.
This proves
$r(\mathcal D_{\mathbf C^*}(X))\subseteq X^{\divis}$.

\emph{Realization on a model of $X$.}
Conversely, write $v=(k/m)\ord_F$, with coprime positive
integers $k,m$. We may take $F$ smooth on a smooth projective
modification $Y\to X$. On $Y\times\mathbf P^1$ use the snc
pair $F\times\mathbf P^1+Y\times\{0\}$ and the monomial
valuation of weights $(k,m)$ along $F\times\{0\}$.
The construction in \cref{prop:mp-monomial-globalization}
realizes this primitive integer valuation as $\ord_E$ after
finitely many ordinary blowups. It is invariant under the action
on $t$, and every center lies over $t=0$, since $\ord_E(t)=m>0$.
If $u$ is a local equation of $F$, then
\[
                 \ord_E(u)=k,\qquad \ord_E(t)=m.
\]
For every ideal on $X$, the normalized restriction is consequently
$(k/m)\ord_F$.

The resulting space is initially a model of $Y$. Its composite
map to $X\times\mathbf P^1$ need not be an isomorphism away
from $t=0$. To obtain a model of $X$, start instead on
$X\times\mathbf P^1$ and successively blow up the global
centers of this same divisorial valuation, normalizing after
each blowup. Apply the extraction argument at a point of its
center on $X\times\mathbf P^1$. Here the local valuation is
already divisorial, so no restriction lemma is needed. It gives
a global prime divisor $D$ with the same primitive order as $E$.

Every center in this second construction is invariant and lies
in the central fiber. Thus each map is an isomorphism over
$\mathbf P^1\setminus\{0\}$. The action lifts to blowups and
normalizations. The resulting normal compact space is K\"ahler,
since the construction is projective over $X\times\mathbf P^1$.
It is flat over $\mathbf P^1$: its local rings are torsion-free
modules over the local DVR of the base. Equivariant resolution,
supported over $0$, gives an snc model if required and preserves
the divisor order. Since $\ord_D(t)=\ord_E(t)=m$, its normalized
restriction is exactly $v$. This is the realization step of
\cite[Proposition~2.5.11]{MPTropicalV2}. Finally the product
model gives $v_{\triv}$, also when $\dim X=0$.
\end{proof}

\begin{example}\label{ex:mp-vertical-normalization-model}
Let $X=\mathbf P^1$ and let $u$ be a parameter at $p\in X$.
The monomial order of weights $(2,3)$ at $(p,0)$ is realized
by a prime divisor $E$ on a model of $X$. Here
\[
 b_E=3,\qquad v_E(t)=1,\qquad
 r(v_E)=\frac23\ord_p,\qquad r(v_E)(\mathcal I_p^3)=2.
\]
Locally the ideal $(u^3,t^2)$ has order six along $E$, and
normalized value two. At a nonzero point of the base the ideal
is the unit ideal. In contrast, passing to $t=0$ sends $t$ to
zero and cannot define the normalized vertical valuation.

The need for the last extraction is visible with
$Y=\operatorname{Bl}_p\mathbf P^2$. The product map
$Y\times\mathbf P^1\to\mathbf P^2\times\mathbf P^1$ restricts
to $\operatorname{Bl}_p\mathbf P^2\to\mathbf P^2$ on every
noncentral fiber, so it is not a model of $\mathbf P^2$ in the
sense of \eqref{eq:mp-analytic-model}. Blowing up the valuation's
centers directly on $\mathbf P^2\times\mathbf P^1$, all over
$0$, is what supplies the required model.
\end{example}
We next identify PL functions with vertical Cartier divisors on
models. Put $T=X\times\mathbf P^1$ and $T_0=X\times\{0\}$.
A \emph{fractional flag ideal} is a $\mathbf C^*$-invariant
coherent $\cO_T$-submodule $\mathfrak b$ of meromorphic
functions satisfying
\[
 \cO_T(-aT_0)\subseteq\mathfrak b\subseteq\cO_T(aT_0)
 \qquad\text{for some integer }a\geq0.
\]
In particular, it is the unit ideal off $T_0$. Near $T_0$ it
admits a finite expression
\[
 \mathfrak b=\sum_j\mathfrak a_jt^{\lambda_j},\qquad
 \lambda_j\in\mathbf Z,\qquad \sum_j\mathfrak a_j=\cO_X.
\]
Here the ideals on $X$ are pulled back to $T$. The expression
is glued to $\cO_T$ off $T_0$: globally, a power $t^q$ denotes
the divisor ideal $\cO_T(-qT_0)$, so it introduces no divisor
at infinity.

To obtain the expression, first twist by a power of $T_0$ so
that $\mathfrak b$ is integral and contains $(t^M)$. Modulo
$t^M$, the action splits the finite $\cO_X$-module
$\bigoplus_{i=0}^{M-1}\cO_Xt^i$ into its weight spaces.
Each weight projection preserves the invariant submodule
$\mathfrak b/(t^M)$: it can be written as a finite linear
combination of the actions of distinct elements of
$\mathbf C^*$ by inverting a Vandermonde matrix. Thus
$\mathfrak b$ has coherent coefficient ideals
$\mathfrak a_i\subseteq\mathfrak a_{i+1}$, with
$\mathfrak a_M=\cO_X$. Undoing the twist gives the asserted
form. This is the flag description in
\cite[equation~(0.0.6)]{MPTropicalV2}.

\begin{lemma}\label{lem:mp-flag-evaluation}
For $v\in X^{\bethan}$, the formulas
\[
 \sigma(v)(\mathfrak b)
    :=\min_j\{v(\mathfrak a_j)+\lambda_j\},\qquad
 \phi_{\mathfrak b}(v):=-\sigma(v)(\mathfrak b)
\]
are independent of the flag expression and take finite real
values. They satisfy
\[
 \begin{aligned}
 \sigma(v)(\mathfrak b\mathfrak c)
     &=\sigma(v)(\mathfrak b)+\sigma(v)(\mathfrak c),\\
 \sigma(v)(\mathfrak b+\mathfrak c)
     &=\min\{\sigma(v)(\mathfrak b),\sigma(v)(\mathfrak c)\},\\
 \sigma(v)(\overline{\mathfrak b})&=\sigma(v)(\mathfrak b).
 \end{aligned}
\]
Moreover, every $\phi\in\PL(X^{\bethan})$ has the form
$m^{-1}(\phi_{\mathfrak b}-\phi_{\mathfrak c})$ for fractional
flag ideals $\mathfrak b,\mathfrak c$ and an integer $m\geq1$.
\end{lemma}

\begin{proof}
For any finite expression, the canonical coefficient ideal
at weight $\lambda$ is
$\sum_{\lambda_j\leq\lambda}\mathfrak a_j$.
Replacing the expression by these increasing coefficient
ideals leaves the minimum unchanged: each original term
occurs at its own weight, and its later occurrences have
larger weight. This proves independence. Since
$\min_jv(\mathfrak a_j)=v(\cO_X)=0$, at least one term is
finite, and the finitely many integer weights give a lower
bound. The identities for products and sums follow by
expanding their coefficient ideals and applying the two
character identities on $X$. This is the Gauss evaluation
on flags of \cite[equation~(2.4.2)]{MPTropicalV2}.

For an integral flag ideal, its integral closure is again an
invariant coherent ideal and contains the same power of $t$.
The reduction argument in the proof of
\cref{prop:mp-integral-closure-values}, now on the compact
space $T$, gives a uniform integer $N$ with
\[
 \overline{\mathfrak b}^{\,N+1}
       =\mathfrak b\,\overline{\mathfrak b}^{\,N}.
\]
All these ideals are flags. Their Gauss values are finite,
so the product identity permits cancellation and proves
the integral-closure assertion. For a fractional flag, twist
by $\cO_T(-aT_0)$ to make it integral. Integral closure
commutes with this invertible twist, and both Gauss values
increase by $a$. This defines the closure independently of
the chosen twist and proves the same assertion.

The finiteness condition in \eqref{eq:mp-pl-plus} is equivalent
to $\sum_j\mathfrak a_j=\cO_X$. Indeed, a common zero $x$
gives the point character that is zero on ideals with unit
stalk at $x$ and infinite on all other ideals; the displayed
maximum would then be $-\infty$. Conversely, a unit sum
has minimum ideal value zero at every character. Hence
\eqref{eq:mp-pl-plus} is $m^{-1}\phi_{\mathfrak b}$ for
$\mathfrak b=\sum_j\mathfrak a_jt^{-k_j}$, and every
$\phi_{\mathfrak b}$ belongs to $\PL^+$. Clearing rational
denominators and using the product identity combines all
positive terms into one flag and all negative terms into
another; an empty product is the unit flag. This proves
the final assertion; compare
\cite[Proposition~2.4.3]{MPTropicalV2}.
\end{proof}

For a model $\cX$, write
$\operatorname{VCar}(\cX)_{\mathbf Q}$ for the space of
$\mathbf C^*$-invariant vertical $\mathbf Q$-Cartier divisors.
The transition maps between these spaces are pullbacks.

\par\noindent\begin{minipage}{\linewidth}
\begin{proposition}\label{prop:mp-pl-vcar}
Evaluation on normalized vertical components induces an isomorphism
\[
 \PL(X^{\bethan})
      \simeq
 \varinjlim_{\cX}\operatorname{VCar}(\cX)_{\mathbf Q}.
\tag{59.14d}\label{eq:mp-pl-vcar}
\]
If $D=\sum_E d_EE$ represents $\phi$ on $\cX$, then
\[
                         d_E=b_E\phi\bigl(r(v_E)\bigr).
\tag{59.14e}\label{eq:mp-vcar-coefficients}
\]
\end{proposition}
\end{minipage}\par

\begin{proof}
\emph{From flags to divisors.}
For $\phi=m^{-1}\phi_{\mathfrak b}$, twist
$\mathfrak b$ to an integral flag and take its normalized
blowup. The fractional pullback is invertible; write
\[
 \mathfrak b\cO_{\cX}=\cO_{\cX}(B_{\mathfrak b}),
 \qquad D=\frac1m B_{\mathfrak b}.
\]
The divisor $B_{\mathfrak b}$ is vertical and invariant,
and is anti-effective when $\mathfrak b$ is integral.
Further domination, including passage to an snc model,
pulls back this equality of fractional ideals. For each
vertical prime $E$, the orders of the flag terms give
\[
 \begin{aligned}
 \sigma(r(v_E))(\mathfrak b)
   &=\min_j\{r(v_E)(\mathfrak a_j)+\lambda_j\}\\
   &=b_E^{-1}\ord_E(\mathfrak b\cO_{\cX}).
 \end{aligned}
\]
Consequently
$b_E\phi(r(v_E))=-m^{-1}\ord_E(\mathfrak b\cO_{\cX})=d_E$.
This proves \eqref{eq:mp-vcar-coefficients} with its sign and
denominator, using the projection pullback in
\eqref{eq:mp-vertical-restriction}.

Products of flags add their divisors. On a common model,
any rational relation among the functions gives the same
relation among every vertical coefficient by
\eqref{eq:mp-vcar-coefficients}. The corresponding divisor
is therefore zero. Thus the assignment descends to a
well-defined linear map from the space of functions $\PL$,
and the pullback equality makes its target the direct limit.

\emph{Injectivity.}
Write $\phi=m^{-1}(\phi_{\mathfrak b}-\phi_{\mathfrak c})$
as in \cref{lem:mp-flag-evaluation}. If its divisor class is
zero, then on a common normal dominating model the
invertible fractional ideals
$\mathfrak b\cO_{\cX}$ and $\mathfrak c\cO_{\cX}$ are equal.
Twist both by the same power of $T_0$ to make them integral.
The proper-modification form of
\eqref{eq:mp-normalized-blowup} shows that their integral
closures on $T$ are equal. By
\cref{lem:mp-flag-evaluation}, their Gauss values agree at
every $v\in X^{\bethan}$. Hence $\phi=0$. This argument
includes proper-support characters and does not use
divisorial density.

\emph{Relative generation and surjectivity.}
Let $D$ be a vertical Cartier divisor; the rational case
follows by clearing a denominator. Pass to a dominating
model $\mu:\cX\to T$ obtained by finitely many equivariant
blowups with centers over $T_0$. Such models are cofinal
by \cite[Remark~2.5.2 and the proof of Theorem~2.6.4,
footnote~10]{MPTropicalV2}.
There is a vertical Cartier divisor $G$ with
$\cO_{\cX}(G)$ relatively ample over $T$. For a single
blowup with center ideal $\mathfrak j$, take the divisor
defined by $\mathfrak j\cO_{\cX}=\cO_{\cX}(G)$; this
is the tautological relatively ample sheaf. For a composite
$f:\cX\to\cY\to T$, combine its divisor $G_f$
with a previously constructed divisor $G_{\cY}$ by
$G=G_f+a f^*G_{\cY}$ for sufficiently large $a$.
The composition rule for relative ampleness is
\cite[Tag~0C4K]{stacks-project}. It applies to the projective
analytic germs through
\cref{thm:mp-bingener-projective-germs}. The finite generation
argument for this rule gives a bound on a neighborhood of
each base point, and a finite cover of $T$ gives a common
bound at each stage. Every divisor used is vertical. For
the identity model one can take $G=0$.

After replacing $G$ by a positive multiple, we may also
assume that $\cO_{\cX}(G)$ is relatively globally generated.
Indeed, on each projective germ a sufficiently large power
is generated, by \cite[Tag~01Q3, assertion~(7)]{stacks-project}
and coherent algebraization in
\cref{thm:mp-bingener-projective-germs}. A finite set of
generating sections spreads to a neighborhood of the base
point; properness removes the support of its cokernel after
shrinking. Choose finitely many such neighborhoods covering
$T$ and a common multiple of their exponents. For
$\cO_{\cX}(D)$ the same argument gives an exponent on each
of finitely many neighborhoods. Taking the largest of
these exponents and tensoring by the already generated
powers of $\cO_{\cX}(G)$ yields a single $N\geq1$ such that
both $\cO_{\cX}(NG)$ and $\cO_{\cX}(D+NG)$ are relatively
globally generated.

Set, as subsheaves of meromorphic functions on $T$,
\[
 \mathfrak b:=\mu_*\cO_{\cX}(D+NG),\qquad
 \mathfrak c:=\mu_*\cO_{\cX}(NG).
\]
These are coherent invariant fractional flag ideals.
To check the bounds in their definition, bound the
coefficients of each vertical divisor between those of
$-a\mu^*T_0$ and $a\mu^*T_0$. This is possible since
all central multiplicities are positive. Push forward
the resulting inclusions, using
$\mu_*\cO_{\cX}=\cO_T$ and the projection formula.
Off $T_0$ both pushforwards are the unit ideal. Relative
generation gives equalities of fractional ideals
\[
 \mathfrak b\cO_{\cX}=\cO_{\cX}(D+NG),\qquad
 \mathfrak c\cO_{\cX}=\cO_{\cX}(NG).
\]
The function $\phi_{\mathfrak b}-\phi_{\mathfrak c}$
therefore maps to $D$. This proves surjectivity, and
injectivity makes the inverse independent of all choices.
The correspondence is \cite[Theorem~2.6.4]{MPTropicalV2}.
\end{proof}

\begin{example}\label{ex:mp-pl-flag-denominator}
Let $X=\mathbf P^1$, fix $p\in X$, and blow up
$(p,0)\in T$. Write $E$ for the exceptional divisor and
$H$ for the strict transform of $T_0$. Then
$\cX_0=H+E$. The flag
$\mathfrak a=\mathcal I_p+(t)$ pulls back to
$\cO_{\cX}(-E)$, so
\[
 \phi(v)=\frac12\max\{\log|\mathcal I_p|(v),-1\}
        =-\frac12\min\{v(\mathcal I_p),1\}
 \quad\longleftrightarrow\quad -\frac12E.
\]
Both central multiplicities are one. Their restrictions are
$r(v_H)=v_{\triv}$ and $r(v_E)=\ord_p$, and the function
has values $0$ and $-1/2$, respectively. At the point
character supported on $p$, the value of $\mathcal I_p$
is infinite but $\phi$ still equals $-1/2$.

Adding $3/2$ replaces the flag by
$\mathfrak b=t^{-3}\mathfrak a$ near $T_0$, glued to
the unit ideal away from $T_0$. Thus
\[
 \phi+\frac32
   =\frac12\max\{\log|\mathcal I_p|+3,2\}
 \quad\longleftrightarrow\quad
 \frac12(3\cX_0-E)=\frac32H+E.
\]
This calculation retains the rational denominator and the
sign in $\mathfrak b\cO_{\cX}=\cO_{\cX}(B_{\mathfrak b})$.
The shift changes only the central fiber, including its
strict transform.
\end{example}
\begin{theorem}\label{thm:mp-divisorial-density}
The rational divisorial locus $X^{\divis}$ is dense in
$X^{\bethan}$.
\end{theorem}

\begin{proof}
Let $F$ be the closure of $X^{\divis}$.  Suppose that
$v_0\notin F$.  The lattice $\PL(X^{\bethan})$ contains the rational constants
and separates points.  For each $w\in F$, separation and the rational
constants give $f_w\in\PL(X^{\bethan})$ with
$f_w(v_0)>0$ and $f_w(w)<0$: start with a separating function and
subtract a rational number strictly between its two values, changing the
sign if necessary.  The sets $\{f_w<0\}$ cover $F$, so
compactness gives $w_1,\ldots,w_N$ for which
\[
 f:=\min_j f_{w_j}\quad\text{satisfies}\quad
 f(v_0)>0,\qquad f|_F<0.
\]
Consequently $\phi:=\max\{f,0\}$ belongs to $\PL$, vanishes on
$F$, and satisfies $\phi(v_0)>0$.

Represent $\phi$ by a vertical divisor $D$ on a single model using
\Cref{prop:mp-pl-vcar}.  For every component $E$ of its central fiber,
\Cref{thm:mp-invariant-divisorial-comparison} gives
$r(v_E)\in X^{\divis}$.  Hence
$b_E\phi(r(v_E))=0$, and
\eqref{eq:mp-vcar-coefficients} gives $D=0$.  The isomorphism
\eqref{eq:mp-pl-vcar} then gives $\phi=0$, contradicting
$\phi(v_0)>0$.  Thus $F=X^{\bethan}$.  This is the proof of
\cite[Theorem~2.3.6, pp.~18 and~28]{MPTropicalV2}; the preceding argument spells out the
point-versus-closed-set lattice separation used there.
\end{proof}

Fix an snc model $\cX$, with central fiber
$\cX_0=\sum_{i\in I}b_iE_i$ as in
\eqref{eq:mp-central-fiber}. For every nonempty $J\subseteq I$
and every connected component $Z$ of
$E_J:=\bigcap_{i\in J}E_i$, take a simplex
\[
 \sigma_Z:=\left\{w\in(\mathbf R_{\geq0})^J\ \middle|\
                     \sum_{i\in J}b_iw_i=1\right\}.
\tag{59.15}\label{eq:mp-dual-simplex}
\]
The snc condition makes $Z$ smooth and irreducible. If
$\varnothing\ne K\subseteq J$, glue the face with weights zero
outside $K$ to the simplex of the unique component of $E_K$
containing $Z$. The resulting finite complex is $\Delta_{\cX}$.
We retain the label $Z$: different components of $E_J$ give
different simplices, even when their vertex sets agree. The
vertex $\xi_i$ corresponding to $E_i$ has weight $1/b_i$, not
$1$. The equality in \eqref{eq:mp-dual-simplex} says that the
monomial valuation has value $1$ on $t$.

Restriction of these monomial valuations along
$\cX\to X\times\mathbf P^1\to X$ gives a continuous map
$i_{\cX}$ by \cref{prop:mp-monomial-globalization}, applied
over $X\times\mathbf P^1$. We shall construct the second map in
\[
 i_{\cX}:\Delta_{\cX}\hookrightarrow X^{\bethan},
 \qquad
 p_{\cX}:X^{\bethan}\longrightarrow\Delta_{\cX},
\tag{59.16}\label{eq:mp-inclusion-retraction}
\]
and prove that it is a continuous left inverse of the first.
For characters of proper support, evaluation on an ideal of
$\cX$ requires a definition; such an ideal is not an ideal on $X$.

Write $\varphi_D\in\PL(X^{\bethan})$ for the function
corresponding to a vertical rational Cartier divisor $D$ by
\cref{prop:mp-pl-vcar}. Let $\mathfrak a\subseteq\cO_{\cX}$
be a coherent invariant ideal equal to the unit ideal off
$\cX_0$. On its normalized blowup, followed by an snc
domination, write
\[
 \mathfrak a\cO_{\cY}=\cO_{\cY}(-F_{\mathfrak a}),
 \qquad u_{\mathfrak a}:=\varphi_{F_{\mathfrak a}}.
\]
Here $F_{\mathfrak a}$ is effective and vertical. Further
domination pulls it back, so the definition is independent
of the chosen domination. If a prime $G$ on a common model
represents $v=r(v_G)\in X^{\divis}$, the coefficient formula
\eqref{eq:mp-vcar-coefficients} gives
$u_{\mathfrak a}(v)=b_G^{-1}\ord_G(\mathfrak a\cO)$.
Thus $u_{\mathfrak a}$ continuously extends the normalized
ideal order. It is nonnegative by divisorial
density. In particular,
$u_{\cO_{\cX}(-E_i)}=\varphi_{E_i}$ and
$\varphi_{\cX_0}=1$. The sign here is positive: the function
associated to the ideal order uses $F_{\mathfrak a}$, whereas
the logarithm of the flag ideal uses $-F_{\mathfrak a}$.

\begin{lemma}\label{lem:mp-retraction}
For every snc model $\cX$, the map $p_{\cX}$ in
\eqref{eq:mp-inclusion-retraction} is well defined and continuous,
and $p_{\cX}i_{\cX}=\mathrm{id}_{\Delta_{\cX}}$.
If $\mu:\cX'\to\cX$ dominates snc models, pullback gives a
continuous facewise integral-affine map
\[
 r_{\cX',\cX}:\Delta_{\cX'}\longrightarrow\Delta_{\cX},
 \qquad r_{\cX',\cX}p_{\cX'}=p_{\cX}.
\tag{59.17}\label{eq:mp-dual-transition}
\]
These transition maps compose under further domination.
\end{lemma}

\begin{proof}
\emph{Strata and continuity.}
Let $\mathcal S$ be the finite set of closed strata $Z$
used in constructing $\Delta_{\cX}$. Each ideal
$\mathcal I_Z$ is invariant and is the unit ideal off
$\cX_0$. Define $h_Z$ on the complex by taking the
monomial value of $\mathcal I_Z$. On a closed simplex
$\sigma_W$ its formula is
\[
 h_Z(w)=
 \begin{cases}
  \min_{i\in J_Z}w_i,&W\subseteq Z,\\
  0,&W\not\subseteq Z,
 \end{cases}
\]
where $Z$ is a component of $E_{J_Z}$. Indeed, along
$W\subseteq Z$ the ideal is generated by the corresponding
boundary equations; otherwise it is the unit ideal at a
general point of $W$. The formulas agree on faces. When a
face lies in a larger stratum not contained in $Z$, at least
one of the coordinates indexed by $J_Z$ has become zero.
Consequently every $h_Z$ is continuous.

These functions separate points of the complex. The functions
$h_{E_i}$ recover the weights, extended by zero to $I$. If
two points have the same weights with positive support $J$
but lie on different components $W,W'$ of $E_J$, then
\[
 h_W(w)=\min_{i\in J}w_i>0,\qquad h_W(w')=0.
\]
Thus $H:=(h_Z)_{Z\in\mathcal S}$ is a continuous injection
of the compact finite complex into $\mathbf R^{\mathcal S}$.
It is a homeomorphism onto its closed image; in particular
$\Delta_{\cX}$ is Hausdorff.

Now consider the continuous map
\[
 A:X^{\bethan}\longrightarrow\mathbf R^{\mathcal S},
 \qquad A(v)=(u_{\mathcal I_Z}(v))_{Z\in\mathcal S}.
\]
For $v\in X^{\divis}$ choose a representing vertical prime
$G$ on a model dominating $\cX$, using
\cref{thm:mp-invariant-divisorial-comparison}. Its center
$C$ on $\cX$ is an irreducible analytic subvariety. Put
$J=\{i:C\subseteq E_i\}$ and
$w_i=b_G^{-1}\ord_G(E_i)$. These weights are positive
exactly for $i\in J$, and
$\sum_i b_iw_i=b_G^{-1}\ord_G(t)=1$.
The unique component $W$ of $E_J$ containing $C$ specifies
a point $w\in\sigma_W$. One has $C\subseteq Z$ exactly
when $W\subseteq Z$: containment specifies both the boundary
indices and the component of their intersection.
The local generators of every $\mathcal I_Z$ therefore give
$u_{\mathcal I_Z}(v)=h_Z(w)$. Hence $A(v)\in H(\Delta_{\cX})$.
Divisorial density and closedness of this image give the same
inclusion for all $v\in X^{\bethan}$. Define
$p_{\cX}:=H^{-1}A$. This proves existence and continuity,
including proper-support characters. In particular,
$w_i(p_{\cX}(v))=\varphi_{E_i}(v)$ for every $i\in I$.

\emph{The section.}
For a rational $w\in\Delta_{\cX}$, the blowup construction in
\cref{prop:mp-monomial-globalization} realizes its monomial
valuation by a vertical prime $G$ on a dominating snc model.
Each blowup center is a component of an intersection of
invariant boundary divisors, so it is invariant under the
connected group $\mathbf C^*$. The resulting prime is invariant
as well. Since the monomial value of $t$ is $1$, the scale
is $b_G^{-1}$. Thus $i_{\cX}(w)=r(v_G)\in X^{\divis}$,
and the preceding construction returns both its weights and
its stratum: $p_{\cX}i_{\cX}(w)=w$.
Rational points are dense in every simplex. The continuity
already proved therefore gives $p_{\cX}i_{\cX}=\mathrm{id}$
everywhere. In particular $i_{\cX}$ is an embedding and
$p_{\cX}$ is surjective.

\emph{Domination.}
Write
\[
 \mu^*E_i=\sum_j a_{ij}E'_j,\qquad
 a_{ij}\in\mathbf Z_{\geq0},\qquad
 b'_j=\sum_i b_i a_{ij}.
\]
On the relative interior of the simplex indexed by $W'$,
send $w'$ to the weights
$w_i=\sum_j a_{ij}w'_j$, in the simplex of the smallest
closed stratum containing $\mu(W')$. A weight $w_i$ is
positive exactly when $\mu(W')\subseteq E_i$: the pullback
divisor has only central components, and at a general point
of $W'$ its local equation is a unit times the corresponding
boundary monomial. The displayed identity for $b'_j$ gives
$\sum_i b_iw_i=1$. On a face the source stratum becomes
larger; its image selects the unique containing component
of the smaller target intersection. Thus the formulas glue
to a continuous facewise integral-affine map $r_{\cX',\cX}$.
Pullback of divisors and containment of these strata also
show that the maps compose.

For a divisorial $v$ represented on a common dominating
model, orders of the pullbacks give the same weight formula.
The center on $\cX$ is the image of the center on $\cX'$.
Its smallest containing stratum agrees with the one just
specified: the boundary divisors containing the image are
exactly those with positive pullback order, and containment
fixes the intersection component. This proves
$r_{\cX',\cX}p_{\cX'}(v)=p_{\cX}(v)$ on $X^{\divis}$.
Both maps are continuous, so density gives the identity on
$X^{\bethan}$. These constructions supply the ideal evaluation
and stratum labels in \cite[Section~3.1.1, pp.~28--29]{MPTropicalV2}.
\end{proof}

For the product model, the complex has a single vertex,
whose image is $v_{\triv}$; the retraction is constant.
This also covers the case $\dim X=0$.

\begin{example}\label{ex:mp-dual-repeated-edges}
Let $X=\mathbf P^1\times\mathbf P^1$. The graphs $C_1$ and
$C_2$ of $z\mapsto z$ and $z\mapsto-z$ meet transversely at
$p=(0,0)$ and $q=(\infty,\infty)$. Blow up $C_1\times\{0\}$
in $X\times\mathbf P^1$, and then blow up the strict transform
of $C_2\times\{0\}$. Denote the final strict transforms of
the original central fiber and the first exceptional divisor
by $H$ and $E$, and the second exceptional divisor by $F$.
Then $\cX_0=H+E+F$ is snc. Near either intersection point,
after the first blowup one has $t=uv$, with the second center
$(u,z)$. Its two blowup charts have
$t=zsv$ and $t=uv$, respectively, which verify the assertion
and the multiplicities.

The intersection $E\cap F$ consists of two disjoint
$\mathbf P^1$ fibers $Z_p,Z_q$, one over each intersection
point. Thus the dual complex has two different edges with
vertices $E,F$. Their midpoints $w_p,w_q$ have the same
component weights $(w_H,w_E,w_F)=(0,1/2,1/2)$, but satisfy
$h_{Z_p}(w_p)=1/2$ and $h_{Z_p}(w_q)=0$.
\end{example}

Write $\operatorname{Aff}_{\mathbf Q}(\Delta_{\cX})$ for the
continuous functions that are affine on each of its simplices
and rational at the vertices. This refers to the given
simplices, before any subdivision.

\begin{theorem}\label{thm:mp-inverse-limit-dual}
The canonical map
\[
 X^{\bethan}\longrightarrow\varprojlim_{\cX}\Delta_{\cX},
 \qquad v\longmapsto(p_{\cX}(v))_{\cX},
\tag{59.18}\label{eq:mp-inverse-limit-dual}
\]
is a homeomorphism. Moreover,
\[
 \PL(X^{\bethan})
   \simeq\varinjlim_{\cX}\operatorname{Aff}_{\mathbf Q}(\Delta_{\cX})
   \simeq\varinjlim_{\cX}\operatorname{VCar}(\cX)_{\mathbf Q}.
\tag{59.19}\label{eq:mp-pl-vcar-dual}
\]
\end{theorem}

\begin{proof}
A vertical divisor $D=\sum_i d_iE_i$ on $\cX$ defines
$f_D(w)=\sum_i d_iw_i$, affine on each simplex, with
$f_D(\xi_i)=d_i/b_i$. Conversely, an element $f$ of
$\operatorname{Aff}_{\mathbf Q}(\Delta_{\cX})$ is determined
on each simplex by its vertex values and corresponds to
$\sum_i b_if(\xi_i)E_i$. This remains true when two
simplices have the same vertices. Pullback of divisors
agrees with pullback by $r_{\cX',\cX}$. Taking direct limits
and applying \cref{prop:mp-pl-vcar} proves
\eqref{eq:mp-pl-vcar-dual}. The construction of $p_{\cX}$
also gives the equality of functions on all of $X^{\bethan}$,
\[
 \varphi_D(v)=\sum_i d_i\varphi_{E_i}(v)
             =f_D(p_{\cX}(v)).
\]
This checks the evaluation identity of
\cite[Lemma~3.1.2, equations~(3.1.2)--(3.1.4)]{MPTropicalV2}
with the normalized vertex coordinates retained.

The canonical map is well defined and continuous by
\cref{lem:mp-retraction}. It is injective because the
functions in $\PL(X^{\bethan})$ separate points and each
factors through one of the maps $p_{\cX}$.
For surjectivity, fix a compatible tuple
$(w_{\cX})_{\cX}$. The sets
\[
                  F_{\cX}:=p_{\cX}^{-1}(w_{\cX})
\]
are closed and nonempty, since $p_{\cX}$ has a section.
For finitely many models choose a common dominating snc
model $\cY$. Compatibility puts $i_{\cY}(w_{\cY})$ in
all the corresponding sets $F_{\cX}$. They therefore have
the finite intersection property. Compactness of $X^{\bethan}$
gives a point in their total intersection, proving surjectivity.
The inverse limit is Hausdorff, so this continuous bijection
from a compact space is a homeomorphism. This proves
\cite[Theorem~3.1.1, pp.~29--31]{MPTropicalV2} using the
retraction sections for the surjectivity step.
\end{proof}

Put $R_{\cX}:=i_{\cX}p_{\cX}$. For a dominating model
$\cY\to\cX$, \cref{lem:mp-retraction} gives
$p_{\cX}i_{\cY}=r_{\cY,\cX}$ and hence
$R_{\cX}R_{\cY}=R_{\cX}$. Moreover, for every
$v\in X^{\bethan}$,
\[
                            R_{\cX}(v)\longrightarrow v
\]
as the snc models are ordered by domination. Indeed, if
$\phi=f_D\circ p_{\cY}\in\PL(X^{\bethan})$, then
$\phi(R_{\cX}(v))=\phi(v)$ for every $\cX$ dominating
$\cY$. Uniform density of $\PL$ in
$C^0(X^{\bethan},\mathbf R)$, from \cref{prop:mp-pl-dense},
gives convergence of every continuous function on this net,
and hence convergence in the compact Hausdorff space
$X^{\bethan}$. No countable cofinal family of models is used.

When \(X\) is projective, this description specializes to the trivially
valued algebraic framework of \cite[Appendix~A]{BJ22}: GAGA identifies the
coherent ideals and projective modifications on the two sides, and the
normalized equation \(\sum b_iw_i=1\) gives the same compact dual complexes.
For a nonprojective compact K\"ahler manifold, the conclusion is analogous
but the proof is not obtained by applying the algebraic theorem to a
meromorphic function field; the ideal spectrum and the Bingener
local-to-global passage replace that unavailable algebraic model.

\section{Log discrepancy}

For a prime divisor \(F\) on a smooth modification \(Y\to X\), put

\[
 A_X(r\ord_F)
   :=r\bigl(1+\ord_F(K_{Y/X})\bigr),
 \qquad r\in\mathbf Q_{>0}.
\tag{59.20}\label{eq:mp-ordinary-discrepancy}
\]

On each monomial cone this extends as the unique linear function with the
prescribed values at the rays.  Under further snc blowups, the pullback
retraction satisfies

\[
              A_X\circ r_{Y',Y}\leq A_X.
\tag{59.21}\label{eq:mp-discrepancy-monotonicity}
\]

For a vertical component \(E\subset\cX_0\), normalized by
\(v_E=b_E^{-1}\ord_E\), define the total discrepancy

\[
 A^{\mathrm{tot}}(v_E)
 :=b_E^{-1}\bigl(1+\ord_E K_{\cX/(X\times\mathbf P^1)}\bigr).
\tag{59.22}\label{eq:mp-total-discrepancy}
\]

It is affine on every dual face.  The function on all of
\(X^{\bethan}\) is

\[
 A^{\mathrm{tot}}(v)
 :=\sup_{\cX}A^{\mathrm{tot}}(p_{\cX}(v))
 \in[1,+\infty].
\tag{59.23}\label{eq:mp-discrepancy-envelope}
\]

The lower bound uses smoothness of \(X\).  The supremum makes
\(A^{\mathrm{tot}}\) lower semicontinuous.  The later entropy convention of
\cite[Definition~2.4.1]{MPWN25} removes the discrepancy of the base divisor
\(X\times\{0\}\):

\[
                    A^{\NA}_X:=A^{\mathrm{tot}}-1.
\tag{59.24}\label{eq:mp-reduced-discrepancy}
\]

\begin{proposition}\label{prop:mp-gauss-discrepancy}
For a divisor \(F\) over \(X\), the Gauss extension \(\sigma(\ord_F)\)
satisfies
\[
 A_{X\times\mathbf P^1}(\sigma(\ord_F))=A_X(\ord_F)+1.
\tag{59.25}\label{eq:mp-gauss-discrepancy}
\]
Thus \(A_X^{\NA}\) agrees with ordinary log discrepancy on Gauss extensions.
\end{proposition}

\begin{proof}
On \(Y\times\mathbf P^1\), the Gauss extension is monomial with weight
\((1,1)\) along \(F\times\mathbf P^1\) and \(Y\times\{0\}\).  Additivity of
monomial discrepancy gives
\[
 A_{X\times\mathbf P^1}(F\times\mathbf P^1)
 +A_{X\times\mathbf P^1}(Y\times\{0\})
 =A_X(F)+1.
\]
This is \cite[(3.2.3)]{MP24}.
\end{proof}

The product model is the normalization unit test.  Its only central
component is \(X\times\{0\}\), with multiplicity one and relative canonical
coefficient zero.  Hence \(A^{\mathrm{tot}}(v_{\triv})=1\) and
\(A_X^{\NA}(v_{\triv})=0\).  After deforming a divisor \(F\) to the normal
cone, the new monomial component represents the Gauss extension of
\(\ord_F\); \eqref{eq:mp-gauss-discrepancy} then reads
\(A_X^{\NA}(\sigma(\ord_F))=A_X(F)\).  These two values detect exactly the
shift that is needed by the entropy convention below.

\begin{warning}{The shift by one}
The total discrepancy \eqref{eq:mp-total-discrepancy} and the reduced
discrepancy \eqref{eq:mp-reduced-discrepancy} are both used in the literature.
The entropy in \cite[Definition~2.4.1]{MPWN25} and in
\Cref{chap:mp-psh,chap:mpwn-cy,chap:mpwn-k-stability} is
\(\int(A^{\mathrm{tot}}-1)\,\rd \mu\).  By contrast,
\cite[(3.2.4) and Definition~6.2.1]{MP24} first constructs and then integrates
the total function.  Omitting the shift changes the Mabuchi functional by the
total mass, which is \(1\) in our normalization.  Thus the reduced convention
is the convention of the later primary source.
\Cref{lem:mp-log-volume-test,thm:mp-reduced-entropy-slope} verify the
corresponding model volume estimate and deduce the reduced entropy lower
bound.  The product-model obstruction to using the total function remains
\eqref{eq:mp-total-slope-obstruction}.
\end{warning}

\section{Compactness, approximation, and nonalgebraic examples}

The preceding results give three distinct approximation mechanisms:

\begin{enumerate}
\item continuous functions are uniformly approximated by PL functions;
\item a point of \(X^{\bethan}\) is approximated by its monomial retractions
      on increasingly fine snc models;
\item monomial points are approximated facewise by rational divisorial points.
\end{enumerate}

The first is functional, the second is birational, and the third is
arithmetic.  They should not be merged into a single sequential assertion.

\begin{example}[A nonprojective complex torus]
Let \(X=\mathbf C^g/\Lambda\), \(g\geq2\), be a very general complex torus.
Then \(X\) is K\"ahler, but \(\NS(X)=0\), so it is not projective and has no
nonzero integral \((1,1)\)-class.  Nevertheless its coherent ideals define
the compact space \(X^{\bethan}\).  Blowing up a point produces an exceptional
divisor \(E\) and the nontrivial valuation \(\ord_E\).  Blowing up invariant
centers in the deformation to the normal cone of the point produces genuine
analytic models and dual complexes.  None of this comes from an ample line
bundle on \(X\).
\end{example}

\begin{example}[A nonprojective K3 surface]
Let \(X\) be a K3 surface with \(\Pic(X)=0\).  Its K\"ahler cone is nonempty,
but every K\"ahler class is transcendental.  Blow up a point
\(p\in X\), and then blow up a point on the first exceptional curve.
The resulting snc exceptional divisor has two components meeting normally,
and its monomial cone gives valuations
\[
 \val_{(w_1,w_2)}\!\left(\sum_{\alpha}c_\alpha z^\alpha\right)
   =\min_{c_\alpha\ne0}(\alpha_1w_1+\alpha_2w_2),
 \qquad w_i\geq0.
\]
Their restrictions to ideals on \(X\) are detected by analytic vanishing
orders.  The construction illustrates why the ideal-theoretic spectrum, not
the meromorphic function field, is the appropriate ambient space.
\end{example}

\section*{Source note}

The tropical spectrum and PL density are MP24 Sections~1.2--1.4.  Analytic
modifications, integral closure, and monomial valuations are Sections~2.1--2.3.
The divisorial-density theorem additionally uses the equivariant model and
vertical b-divisor analysis of Sections~2.4--2.6.  The inverse-limit theorem
and total-discrepancy construction are Sections~3.1--3.2; the projective
algebraic comparison is \cite[Appendix~A]{BJ22}.  Reduced
discrepancy is \cite[Definition~2.4.1]{MPWN25}.  The examples above use
only standard nonprojective K\"ahler manifolds; they are included to emphasize
that the MP construction is genuinely analytic and not an algebraic theory in
disguise.  The Stein-local equivalence is traced to Bingener's Aussagen~5.1
and~5.3 through \cite[Theorem~5.4.4 and
Corollaries~5.4.12--5.4.14]{XiaStein26}; the dimension step is exactly
\cite[Aussage~2.8]{Bingener76Stein}, with properness and finite presentation
retained.
